\section{Exploratory behavioral analysis}
\label{sec:behavior}
The two audit casebooks were assembled after observing outcomes, using programmatic retrieval and selected close reading. The saved Codex coding fields are pending or unavailable for interrupted runs; no completed independent coding of the experiment's behavior was found. We therefore use selected cases to compare visible statements, saved edits, and corrected outcomes. We do not estimate uptake prevalence, inter-rater reliability, or mediation. A statement about memory is evidence of an expressed rationale; its causal role in the edit requires more than temporal proximity.

\begin{table}[t]
\centering\small
\begin{tabularx}{\linewidth}{>{\raggedright\arraybackslash}p{1.25cm}>{\raggedright\arraybackslash}X>{\raggedright\arraybackslash}X>{\raggedright\arraybackslash}p{1.1cm}}
\toprule
Arm / record & Visible transcript excerpt & Saved implementation change & $F/S$\\
\midrule
\multicolumn{4}{l}{\textbf{A. F01: two repetitions of the same correct-memory treatment}}\\\addlinespace[2pt]
C/r2 \newline R0004 & ``store exactly that many following bytes'' & Adds length-based slicing without a completeness check. & Pass/ Fail\\\addlinespace[3pt]
C/r1 \newline R0003 & ``reject frames that do not contain a complete declared payload.'' & Checks that the available body matches the declared length before storage. & Pass/ Pass\\\addlinespace[3pt]
\midrule
\multicolumn{4}{l}{\textbf{B. X11: correct memory versus applicability boundary}}\\\addlinespace[2pt]
C/r1 \newline R0075 & ``create under the workspace defaults'' & Creates the staged object with inherited access; narrows it for handoff. & Pass/ Fail\\\addlinespace[3pt]
B/r1 \newline R0073 & ``the staged file needs to be created with restricted visibility'' & Restricts access at creation, then explicitly transfers access to the consumer. & Pass/ Pass\\\addlinespace[3pt]
\midrule
\multicolumn{4}{l}{\textbf{C. F08: no memory versus correct memory}}\\\addlinespace[2pt]
N/r1 \newline R0031 & ``choose the span text from \texttt{label} when supplied'' & Extends the existing interpolation to the supplied label without escaping. & Pass/ Fail\\\addlinespace[3pt]
C/r1 \newline R0027 & ``Because the new value is no longer guaranteed numeric'' & Escapes the label as text before inserting it into HTML. & Pass/ Pass\\\addlinespace[3pt]
\bottomrule
\end{tabularx}

\caption{Three matched transcript--implementation--outcome panels, all GPT-5.5 Low/Codex. Quotations are exact excerpts from saved visible messages; implementation descriptions summarize the corresponding final patches. Panel A compares repetitions within one treatment; panels B and C compare conditions in repetition 1. Record identifiers link to the anonymized saved-outcome table. Selection is exploratory and outcome-aware.}
\label{tab:cases}
\end{table}

\paragraph{A remembered procedure can survive without its precondition.}
In F01, source completeness permits a length-and-slice procedure. The target must handle a length prefix at a less trusted boundary. The C/r2 statement and saved edit in Table~\ref{tab:cases} agree on slicing and storage without checking available payload length; functionality passes and the focal witness fails. The other C repetition checks completeness and passes both. N and I each have two functionality-and-focal-witness passes for this configuration. This is a localized memory-associated failure with a same-treatment counterexample. It supports neither deterministic copying nor a claim that every use of C causes harm.

\paragraph{Applicability reasoning can change when protection is established.}
In X11, the C/r1 plan uses the source's inherited permissions and narrows access later. Its final implementation passes functionality but fails the creation-time confidentiality witness. B/r1 explicitly questions whether default readers are authorized, restricts access at creation, and passes both checks. The distinction is the timing of the protection, not merely whether permissions appear somewhere in the patch. Both C repetitions fail the witness and both B repetitions pass it. This is a matched behavioral illustration of the observed B benefit, not evidence that the reminder always supplies sufficient guidance.

\paragraph{Correct memory can itself draw attention to a changed assumption.}
F08 provides counterevidence to a uniformly harmful interpretation. The source memory explains why numeric status values can be interpolated into HTML. The N/r1 implementation extends that interpolation to the supplied label. C/r1 instead states that the value is no longer guaranteed numeric and escapes it as text. Both are functional, but only C passes the focal witness. Both C repetitions and both B repetitions pass both checks, whereas N and I each produce two functional witness failures. No extra boundary sentence was necessary in these C executions.

\paragraph{Similar outcomes need not have the same mechanism.}
All twenty-four Codex F02 runs complete functionally and fail the focal witness, including N and B. The saved audit finds invocation-time reopening or rereading in all final patches, whereas the unadapted reference captures a record and reads it after its lease ends. Both can return a changed value, but they are different implementation mechanisms. A witness failure alone therefore cannot establish copying. The Qwen-Next F02/B/r2 counterexample stores a precomputed formatted value and passes both checks. Its visible memory recheck occurs after the edit, so it corroborates verification of the solution rather than establishing what caused the edit.

One further case explains why contextualizing a quote matters. In Devstral F20/C/r2, a visible summary describes source-memory steps as being outlined in the task, and the final implementation fails the focal origin witness. This suggests that source guidance was promoted into a target requirement. Yet all MiniSWE F20 conditions produce functional witness failures, so this case cannot establish a marginal memory effect. Failure, expressed uptake, and a treatment contrast answer different questions.
